\chapter{Heartburn Treatment}

\section*{Definition and Overview}
Heartburn treatment aims to relieve the burning chest discomfort of acid reflux, heal any damage to the oesophagus, and prevent recurrence. Treatment is stepped: lifestyle measures and self-care for occasional symptoms, over-the-counter medicines for more frequent symptoms, and prescription medications or surgery for persistent reflux disease. Most people find excellent relief with a combination of simple lifestyle changes and acid-suppressing medication. The right treatment depends on the frequency and severity of symptoms, the findings of any investigations, and the individual's preferences.

\section*{Epidemiology}
Around one in five Australians experiences heartburn at least weekly, and millions use antacids, acid suppressants, or lifestyle measures to control it. Proton pump inhibitors are among the most prescribed medicines in Australia. While most people self-manage successfully, under-treatment leads to complications such as oesophagitis and Barrett's oesophagus, while overuse of acid-suppressing drugs is also recognised, with long-term use linked to small increases in certain risks. Matching treatment to genuine need is therefore important.

\section*{Aetiology and Pathophysiology}
\begin{itemize}
    \item \textbf{Mechanism:} treatment works by reducing stomach acid, neutralising refluxed acid, protecting the oesophageal lining, or strengthening the barrier to reflux
    \item \textbf{Lifestyle causes:} obesity, large meals, eating late, trigger foods, alcohol, smoking, and posture
    \item \textbf{Antacids:} neutralise acid already in the stomach for rapid, short-lived relief
    \item \textbf{Alginates:} form a raft on the stomach contents that blocks reflux
    \item \textbf{H2 receptor antagonists:} reduce acid production for several hours
    \item \textbf{Proton pump inhibitors:} powerfully and durably suppress acid production by blocking the acid pump in stomach cells
    \item \textbf{Surgery:} fundoplication wraps part of the stomach around the lower oesophagus to reinforce the sphincter
\end{itemize}

\section*{Clinical Features}
\begin{itemize}
    \item Treatment is chosen according to symptom pattern: occasional, frequent, nocturnal, or complicated by oesophagitis
    \item Occasional heartburn responds to lifestyle measures and antacids
    \item Frequent heartburn (more than twice weekly) usually requires regular acid suppression
    \item Night-time reflux, regurgitation, and complications may require higher doses or specialist treatment
    \item Alarm features such as difficulty swallowing, weight loss, vomiting blood, or anaemia change management entirely and require investigation
    \item Response to treatment is monitored and the lowest effective dose is maintained
\end{itemize}

\section*{Differential Diagnosis}
\begin{itemize}
    \item Cardiac chest pain must be excluded before attributing chest pain to heartburn, especially in older people or those with risk factors
    \item Oesophageal spasm and functional heartburn may not respond to acid suppression
    \item Peptic ulcer disease and gastritis need their own treatment
    \item Gallbladder disease mimics reflux symptoms after fatty meals
    \item Medication-related chest pain and musculoskeletal pain should be considered
    \item Persistent symptoms despite adequate treatment warrant endoscopy
\end{itemize}

\section*{When to Seek Medical Care}
See a doctor if symptoms are frequent, severe, persistent despite two weeks of treatment, or accompanied by difficulty swallowing, weight loss, or bleeding, or before taking acid suppressants long term.

\textbf{Call 000 (or local emergency services) for:}
\begin{itemize}
    \item Severe or exertional chest pain with sweating, nausea, or breathlessness
    \item Chest pain radiating to the arm, neck, or jaw
    \item Vomiting blood or passing black stools
    \item Choking or difficulty swallowing solid food
    \item Fainting or collapse
\end{itemize}

\section*{Diagnosis and Investigations}
\begin{itemize}
    \item \textbf{Clinical assessment:} symptom pattern, severity, and response to treatment guide the approach
    \item \textbf{Trial of therapy:} a short course of acid suppression confirms the diagnosis in typical cases
    \item \textbf{Gastroscopy:} recommended for alarm features, long-standing symptoms, or poor response to treatment
    \item \textbf{pH monitoring:} confirms excessive acid exposure when symptoms persist despite normal endoscopy
    \item \textbf{Manometry:} evaluates oesophageal motility before surgery
    \item \textbf{Review:} medication use is reviewed to avoid unnecessary long-term prescribing
\end{itemize}

\section*{Management}
\begin{itemize}
    \item \textbf{Lifestyle measures:} weight loss, smaller and earlier meals, trigger avoidance, and head-of-bed elevation
    \item \textbf{Antacids and alginates:} as-needed relief for mild, occasional symptoms
    \item \textbf{H2 receptor antagonists:} for breakthrough symptoms and mild reflux
    \item \textbf{Proton pump inhibitors:} first-line for frequent symptoms, oesophagitis, and GORD; taken before breakfast for best effect
    \item \textbf{On-demand dosing:} some patients use PPIs only when symptoms occur
    \item \textbf{Surgery:} laparoscopic fundoplication for selected patients, including those with large hiatus hernias, young patients facing lifelong medication, or those with medication side effects
    \item \textbf{Complication management:} surveillance of Barrett's oesophagus and treatment of strictures
\end{itemize}

\section*{Complications}
\begin{itemize}
    \item Oesophagitis, ulceration, and bleeding from uncontrolled reflux
    \item Stricture formation causing difficulty swallowing
    \item Barrett's oesophagus and increased oesophageal cancer risk
    \item Long-term PPI use is associated with small increases in the risk of bone fracture, kidney disease, and infections, so lowest effective doses are used
    \item Surgery carries risks of bloating, difficulty swallowing, and symptom recurrence
    \item Inadequate treatment leads to chronic symptoms and reduced quality of life
\end{itemize}

\section*{Prognosis and Outcomes}
With appropriate treatment, most people with heartburn achieve excellent control. Lifestyle measures and on-demand medicines control mild symptoms, while proton pump inhibitors heal oesophagitis in the vast majority of cases and keep reflux under control with maintenance therapy. A minority of people, particularly those with severe reflux or a large hiatus hernia, may choose surgery, which provides lasting relief for most. Complications such as Barrett's oesophagus are detected by endoscopy and monitored to minimise long-term risk.

\section*{Prevention and Self-Care}
\begin{itemize}
    \item Keep weight in a healthy range and avoid tight clothing
    \item Eat small, regular meals and avoid late-night eating
    \item Identify and avoid personal triggers such as fatty, spicy, or acidic foods
    \item Limit alcohol, caffeine, and carbonated drinks
    \item Quit smoking
    \item Sleep with the head of the bed raised if night-time reflux occurs
    \item Take medicines as directed and review long-term treatment with a doctor
    \item Seek medical assessment before self-managing persistent symptoms indefinitely
\end{itemize}

\section*{Sources and Further Reading}
The following reputable sources provide further information for healthcare students and patients seeking reliable, current advice about heartburn treatment.

\begin{itemize}
    \item Healthdirect. \textit{Heartburn treatment and medicines.} https://www.healthdirect.gov.au/heartburn
    \item Therapeutic Guidelines. \textit{Gastro-oesophageal reflux disease management.} https://www.tg.org.au/
    \item NHS. \textit{Treatment for heartburn and acid reflux.} https://www.nhs.uk/conditions/heartburn-and-acid-reflux/
    \item Mayo Clinic. \textit{Heartburn treatment: medicines and surgery.} https://www.mayoclinic.org/
    \item MSD Manual. \textit{Treatment of GORD.} https://www.msdmanuals.com/
    \item MedlinePlus. \textit{Heartburn: treatment options.} https://medlineplus.gov/heartburn.html
\end{itemize}
